\title{Direct Preference Optimization: Your Language Model is Secretly a Reward Model}

\begin{abstract}
Large-scale unsupervised language models (LMs) possess extensive general knowledge and exhibit certain reasoning capabilities; however, precisely steering their output remains challenging because their training lacks supervision. Contemporary approaches to improving steerability typically involve collecting human judgments on the relative merit of model outputs and then fine-tuning the LM accordingly, commonly via reinforcement learning from human feedback (RLHF). Despite its popularity, RLHF tends to be complicated and sometimes unstable, as it requires first training a reward model to encode human preferences and then optimizing the original LM with reinforcement learning to maximize the surrogate reward—careful to avoid significant divergence from the pretrained model. This work introduces a novel reward model parameterization within the RLHF framework, enabling direct closed-form extraction of the optimal policy. Consequently, we reformulate the standard RLHF optimization as a straightforward classification problem. Our proposed approach, \textit{Direct Preference Optimization} (DPO), is stable, effective, and computationally efficient, removing the requirement for sampling from the LM during adaptation and greatly reducing the necessity for hyperparameter adjustments. Through our experiments, we demonstrate that DPO achieves human preference alignment on par with, or superior to, established methods. In particular, DPO-based fine-tuning outperforms RLHF approaches based on PPO in regulating generative sentiment and attains comparable or enhanced results in both summarization and single-turn dialogue—all while offering greater implementation and training simplicity.
\end{abstract}

\section{Introduction}
Extensive unsupervised training of language models (LMs) on massive textual corpora has resulted in the emergence of impressive and sometimes unexpected competencies~\citep{chowdhery2022palm, brown2020language, touvron2023llama,bubeck2023sparks}. Nevertheless, the training data for these models, being authored by humans with diverse intentions, priorities, and expertise, incorporates a broad array of behaviors—some of which may be undesirable for the model to emulate. For instance, although it is beneficial for an AI-powered programming assistant to \textit{recognize} common programming errors in order to correct them, the generation of code should preferentially reflect the higher-caliber examples—however infrequent—they encountered during training. Similarly, while it is important that a language model be \textit{cognizant} of misconceptions commonly held by, say, 50\% of people, it is clearly unacceptable for the model to propagate such misconceptions at the same rate when asked about the corresponding topic. This highlights the importance of appropriately selecting a model’s \emph{preferred outputs and conduct} from its extensive repertoire of \textit{knowledge and skills} in order to produce AI systems that are safe, reliable, and governable~\citep{ouyang2022training}. Although established techniques tend to guide LMs towards human-aligned behavior by employing reinforcement learning (RL) strategies, our work will demonstrate that the RL-based preference learning objectives can in fact be solved exactly using a straightforward binary cross-entropy loss, thereby streamlining the overall preference learning framework.

In summary, current approaches encode preferred behaviors into language models via carefully assembled human preference datasets that encapsulate the actions regarded as desirable or safe. This stage of preference modeling typically follows an initial phase of large-scale, unsupervised pre-training on broad text collections. While supervised fine-tuning on selected human-authored responses is the most direct route to learning preferences, the prevailing and most effective methodology is reinforcement learning from human or AI-generated feedback (RLHF/RLAIF; \citep{christiano2017deep,bai2022constitutional}). The RLHF paradigm involves fitting a reward model to human preference data, then applying RL to train the language model’s policy to yield responses that the reward model evaluates highly, all while constraining deviations from the base model. Despite yielding models with remarkable conversational and programming performance, RLHF introduces significant complexity relative to supervised learning; the process entails training multiple language models, as well as in-loop sampling from policies during training, resulting in substantial computational expenditure.

In this work, we demonstrate a method for directly aligning language models with human preferences, without depending on explicit reward modeling or reinforcement learning. We introduce \textit{{Direct Preference Optimization} (DPO)}, a procedure that implicitly targets the same optimization objective as established RLHF approaches (i.e., maximizing reward subject to a KL-divergence regularization), yet remains easy to implement and straightforward to train. At a high level, DPO works by amplifying the ratio of the log likelihoods assigned to preferred versus non-preferred outputs, incorporating a flexible, per-sample weighting mechanism to avoid the model collapse that can arise from a naive probability ratio formulation. DPO, akin to prior algorithms, utilizes a formal preference model (for example, the Bradley-Terry model; \cite{bradley1952rankanalysis}) to quantify how well a reward function reflects actual preference judgments. Nevertheless, instead of training a reward model via a preference loss and subsequently optimizing a policy with respect to this learned reward, DPO leverages a variable transformation, expressing the preference loss directly in terms of the policy itself. With a collection of human-labeled preferences on model outputs, DPO is able to directly update the policy via a binary cross-entropy loss, effectively producing the policy optimal for an implicit reward function consistent with observed preferences.

Our central contribution is Direct Preference Optimization (DPO), an uncomplicated approach for preference-driven language model training that does not require reinforcement learning. Empirically, we find that DPO matches or outperforms prior methods, including PPO-based RLHF, in a variety of preference-driven tasks—such as sentiment adjustment, text summarization, and conversational response generation—even with models containing up to 6 billion parameters.

\section{Related Work}

Large-scale self-supervised language models have demonstrated the ability to perform certain tasks without explicit training data—either in a zero-shot manner \citep{radford2019language} or when provided with minimal prompting \citep{gpt3,megatron,chowdhery2022palm}. Nonetheless, their capabilities on downstream benchmarks and their alignment with user goals are substantially enhanced when the models are further adapted—through a process known as fine-tuning—using curated datasets of instructions and human-authored responses \citep{mishra-etal-2022-cross,sanh2022multitask,chung2022scaling,thoppilan2022lamda}. This process, referred to as `instruction-tuning', allows large language models (LLMs) to handle instructions not present in the training set and leads to more generally effective models \citep{chung2022scaling}. While instruction tuning has been effective, gathering \textit{relative} human assessments of response quality is often less labor-intensive than assembling expert demonstrations. Consequently, subsequent research has leveraged datasets that capture human preferences to further fine-tune LLMs, yielding performance improvements across domains such as translation \citep{kreutzer-etal-2018-reliability}, summarization \citep{stiennon2022learning,ziegler2020finetuning}, narrative generation \citep{ziegler2020finetuning}, and compliance with instructions \citep{ouyang2022training,ramamurthy2023is}. Typically, these strategies first involve training a neural reward model that aligns with recorded preferences using a model like the Bradley-Terry framework \citep{bradley1952rankanalysis}. Subsequently, the language model is optimized to maximize this learned reward, most often employing reinforcement learning approaches such as REINFORCE \citep{williams1992reinforce}, proximal policy optimization (PPO; \cite{schulman2017proximal}), or similar techniques \citep{ramamurthy2023is}. Related advancements involve using instruction-following LLMs refined via human feedback to generate new synthetic preference data targeting attributes like safety or non-harmfulness \citep{bai2022constitutional}, relying on limited human oversight provided in the form of rubrics to guide the model's outputs. These developments reflect a synthesis of two research directions: the first focuses on using reinforcement learning to train language models towards specific goals~\citep{Ranzato2015SequenceLT,paulus2018a,wu2018learning}; the second centers on methodologies for training models using human preference signals \citep{christiano2017deep,kupcsik2018learning}. Even with the advantages of leveraging comparative human preferences, reinforcement learning-based fine-tuning of large language models remains technically demanding. This study introduces an approach that is grounded in theory and achieves optimization for relative preferences without depending on reinforcement learning.

Beyond the domain of language, policy learning from preferences has been explored within both the bandit and reinforcement learning paradigms, with numerous methods introduced over time. When contextual bandits utilize preferences or action rankings instead of explicit rewards, the framework is referred to as the contextual dueling bandit (CDB; \cite{yue2012karmed,dudik2015contextual}). Without direct access to absolute rewards, theoretical treatments of CDBs replace the standard concept of an optimal policy with the \textit{von Neumann winner}: a policy achieving at least a 50\% expected win rate against every alternative policy \citep{dudik2015contextual}. A notable distinction in the CDB setup is that preference signals are obtained in an online manner, while learning from human preferences generally relies on a static, offline collection of preference-labeled action pairs \citep{yan2022human}. Likewise, \textit{preference-based RL} (PbRL) relies on binary preference signals sourced from an \textit{unknown} underlying `scoring' function rather than directly observed rewards \citep{BusaFekete2014,ruiz2023dueling}. PbRL encompasses several strategies, including those that support reuse of off-policy preference information; typically, these approaches require explicit modeling of the latent scoring (or reward) function before proceeding to policy optimization \citep{jain2013learning,BusaFekete2014,christiano2017deep,sadigh2017active,kupcsik2018learning}. In contrast, our method proposes a unified approach that bypasses reward estimation and focuses on direct policy optimization according to observed preferences.

\section{Preliminaries}\label{section:prelims}

Here, we briefly recapitulate the RLHF pipeline described in \citeauthor{ziegler2020finetuning}, and subsequently expanded in works such as \citep{stiennon2022learning, bai2022training, ouyang2022training}. This pipeline typically comprises three core stages: (1) supervised fine-tuning (SFT); (2) collection of preference data and training of a reward model; and (3) reinforcement learning-based optimization.

\textbf{SFT}: The RLHF process commences by further tuning a large pre-trained language model on curated, task-relevant datasets (covering domains like dialogue, summarization, etc.) through supervised learning, resulting in the fine-tuned policy $\pi^\text{SFT}$.

\textbf{Reward Modelling Phase}: During the second step, the SFT policy generates responses $(y_1, y_2)\sim \pi^\text{SFT}(y \mid x)$ for prompts $x$. Human annotators are then asked to indicate a preference between these candidate outputs, denoted by $y_w\succ y_l \mid x$, where $y_w$ is the selected answer and $y_l$ is the less preferred one among $(y_1, y_2)$. These preference signals are assumed to originate from an unknown, latent reward function $r^*(y, x)$, to which we have no direct access. Multiple approaches have been proposed for modeling such preferences, with the Bradley-Terry (BT) \cite{bradley1952rankanalysis} model widely utilized (more general schemes like the Plackett-Luce model \citep{plackett1975analysis, luce2012individual} are equally applicable when access to lists of ranked responses is available). In the BT model, the preference distribution expressed by humans, denoted as $p^*$, is formalized as follows:

\begin{equation}\label{eq:bradley-terry}
    p^*(y_1\succ y_2 \mid x)=\frac{\exp\left(r^*(x, y_1)\right)}{\exp\left(r^*(x, y_1)\right) + \exp\left(r^*(x, y_2)\right)}.
\end{equation}

Given access to a fixed dataset of preference comparisons $\mathcal{D}=\bigl\{x^{(i)}, y_w^{(i)}, y_l^{(i)}\bigr\}_{i=1}^N$ that have been drawn from $p^*$, one can introduce a reward model $r_{\phi}(x, y)$ and learn its parameters using maximum likelihood estimation. When reframed as a binary classification task, the training objective takes the form of a negative log-likelihood loss:

\begin{equation}\label{eq:reward_model}
    \mathcal{L}_R(r_{\phi}, \mathcal{D}) = -\mathbb{E}_{(x, y_w, y_l)\sim \mathcal{D}}\bigl[\log \sigma(r_{\phi}(x, y_w)- r_{\phi}(x, y_l))\bigr]
\end{equation}

where $\sigma$ denotes the logistic sigmoid. In large language models, it is standard practice to initialize the network $r_{\phi}(x, y)$ from the SFT policy $\pi^\text{SFT}(y \mid x)$, typically by adding a linear head to the output of the last transformer layer to yield a scalar reward value~\cite{ziegler2020finetuning}. To further reduce variance in the reward estimates, previous studies enforce normalization such that $\mathbb{E}_{x,y\sim \mathcal{D}}\left[r_\phi(x, y)\right] = 0$ for each context $x$.

\textbf{RL Fine-Tuning Phase}: In the reinforcement learning stage, this trained reward function supplies feedback for model updates. In line with earlier work~\citep{jaques2017sequence, jaques2020human}, the optimization objective is formalized as:

\begin{equation}\label{eq:RL}
\max_{\pi_{\theta}}  \mathbb{E}_{x\sim \mathcal{D}, y\sim \pi_{\theta}(y \mid x)}\bigl[r_{\phi}(x, y)\bigr] - \beta\mathbb{D}_{\textrm{KL}}\bigl[\pi_{\theta}(y\mid x)\mid \mid \pi_\text{ref}(y\mid x)\bigr],
\end{equation}

where the parameter $\beta$ regulates how much the policy $\pi_\theta$ can deviate from the reference distribution $\pi_\text{ref}$, which is taken to be the SFT baseline $\pi^\text{SFT}$. The initial parameters of $\pi_\theta$ are likewise set to match those of $\pi^\text{SFT}$. This regularization serves two roles: it maintains the language model's outputs within the regime where the reward function is reliable, and also prevents degenerate solutions such as collapsing onto narrow sets of high-reward responses. Owing to the discrete outputs of language models, the above objective cannot be differentiated directly and therefore is usually handled by reinforcement learning methods. The predominant technique~\cite{ziegler2020finetuning, stiennon2022learning, bai2022training, ouyang2022training} defines the shaped reward as ${r(x, y) = r_{\phi}(x, y) -\beta (\log \pi_{\theta}(y\mid x) - \log \pi_\text{ref}(y\mid x))}$ and applies proximal policy optimization (PPO)~\cite{schulman2017proximal} for policy improvement.

\section{Direct Preference Optimization}\label{sec:DPO}

Because reinforcement learning techniques are often unwieldy when scaling up to large models, particularly in the context of language model fine-tuning, our objective is to develop a straightforward method for directly optimizing policies based on preference data. Contrary to prior RLHF approaches—which require training a reward predictor and then optimizing it via RL—we propose to leverage a reward model parameterization that allows the closed-form extraction of its optimal policy, thereby bypassing the need for a reinforcement learning loop. Central to our method is an analytical mapping from reward functions to their corresponding optimal policies; this enables the reparameterization of reward-based loss functions as direct losses over policies. In this way, we circumvent the separate reward modeling step while still ensuring alignment with established models of human

\textbf{Derivation of the DPO Objective.} We begin from the standard reinforcement learning objective as expressed in Eq.~\ref{eq:RL}, where the reward function $r$ is general. Previous research~\citep{peters2007reinforcement, peng2019advantage, korbak2022reinforcement, go2023aligning} has demonstrated that the KL-regularized reward maximization problem in Eq.~\ref{eq:RL} is optimally solved by a policy of the form:

\begin{equation}\label{eq:op_policy}
    \pi_r(y\mid x) = \frac{1}{Z(x)}\pi_\text{ref}(y\mid x)\exp\left(\frac{1}{\beta}r(x, y)\right),
\end{equation}

where $Z(x) =\sum_{y}\pi_\text{ref}(y\mid x)\exp\left(\frac{1}{\beta}r(x, y)\right)$ serves as the normalizing partition function. A detailed derivation is available in Appendix \ref{app:derivation1}. Notably, even when the MLE approximation $r_{\phi}$ is used in place of the actual reward function $r^*$, it remains computationally difficult to compute the partition function $Z(x)$, as discussed in~\citep{korbak2022reinforcement, go2023aligning}. This computational bottleneck makes the direct utilization of this formulation challenging. To overcome this, we can rearrange Eq.~\ref{eq:op_policy} to represent the reward as a function of the corresponding optimal policy $\pi_r$, the reference distribution $\pi_\text{ref}$, and the partition function $Z(\cdot)$. In particular, applying the logarithm to both sides of Eq.~\ref{eq:op_policy} and isolating $r(x,y)$ gives:

\begin{equation}\label{eq:main_eq}
    r(x,y) =\beta \log \frac{\pi_r(y\mid x)}{\pi_\text{ref}(y\mid x)} + \beta \log Z(x).
\end{equation}

We use this parameterization for the true reward $r^*$ and its associated optimal policy $\pi^*$. Importantly, the Bradley-Terry framework relies exclusively on the reward gap between two candidate outputs, i.e., ${p^*(y_1 \succ y_2 \mid x) = \sigma(r^*(x, y_1) - r^*(x, y_2))}$. When the expression from Eq.~\ref{eq:main_eq} is substituted for $r^*(x,y)$ in the preference model (Eq.~\ref{eq:bradley-terry}), the terms involving the partition function are eliminated. Consequently, the probability that a human prefers $y_1$ to $y_2$ can be expressed solely via the optimal policy $\pi^*$ and the reference policy $\pi_\text{ref}$. Thus, under the Bradley-Terry preference model, the optimal RLHF policy $\pi^*$ adheres to:

\begin{equation}\label{eq:objective}
    p^*(y_1\succ y_2 \mid x)=\frac{1}{1 + \exp\left(\beta \log \frac{\pi^*(y_2\mid x)}{\pi_\text{ref}(y_2\mid x)} - \beta \log \frac{\pi^*(y_1\mid x)}{\pi_\text{ref}(y_1\mid x)}\right)}
\end{equation}

A full derivation of this result is provided in Appendix~\ref{app:derivation2}. Although Eq.~\ref{eq:objective} is presented for the Bradley-Terry model, analogous forms can be derived for broader classes of preference models such as the Plackett-Luce models~\citep{plackett1975analysis, luce2012individual}, as elaborated in Appendix~\ref{app:plackett_luce_models}.

With the human preference likelihood now written directly in terms of the optimal policy, we are able to formulate a likelihood maximization problem for a parameterized policy $\pi_\theta$. Mirroring the approach taken in reward modeling (cf. Eq.~\ref{eq:reward_model}), we obtain the following objective for optimizing the policy:

\begin{equation}\label{eq:optimum_model}
    \mathcal{L}_\text{DPO}(\pi_{\theta}; \pi_\text{ref}) = -\mathbb{E}_{(x, y_w, y_l)\sim \mathcal{

In this formulation, we learn an implicit reward through a distinct parameterization, wherein the optimal policy directly corresponds to $\pi_\theta$. Additionally, because this approach mirrors the process of fitting a reparameterized Bradley-Terry model, it inherits certain theoretical guarantees, including consistency, assuming the preference data are sampled from appropriate distributions \cite{bong2022generalized}. We examine the theoretical underpinnings of DPO and compare them to related frameworks in Section~\ref{sec:theory}.

\textbf{Mechanics of the DPO Update.} To gain insight into how DPO operates, it is helpful to inspect the gradient of the loss function $\mathcal{L}_\text{DPO}$. Specifically, the gradient with respect to $\theta$ takes the following form:

\begin{multline*}\label{eq:gradient}
    \nabla_\theta \mathcal{L}_\text{DPO}(\pi_\theta;\pi_\text{ref}) = \\ -\beta\mathbb{E}_{(x, y_w, y_l) \sim \mathcal{D}} \bigg[\underbrace{\sigma(\hat{r}_\theta(x, y_l) - \hat{r}_\theta (x, y_w))}_\text{higher weight when reward estimate is wrong}\bigg[\underbrace{\nabla_\theta\log \pi(y_w \mid x)}_\text{increase likelihood of $y_w$} - \underbrace{\nabla_\theta\log\pi(y_l \mid x)}_\text{decrease likelihood of $y_l$}\bigg]\bigg],
\end{multline*}

where $\hat{r}_\theta(x, y) = \beta \log \frac{\pi_\theta(y \mid x)}{\pi_\text{ref}(y \mid x)}$ represents the implicit reward conferred by the language model $\pi_\theta$ with respect to the reference model $\pi_\text{ref}$ (further discussed in Section~\ref{sec:theory}). Conceptually, this gradient formulation amplifies the likelihood of preferred completions $y_w$ while diminishing that of less favored ones $y_l$. Crucially, each data instance is weighted according to the extent to which the implicit reward model $\hat{r}_\theta$ erroneously prefers the disfavored completion, with this discrepancy scaled by the factor $\beta$. This weighting effectively quantifies the misranking of completions and incorporates the KL penalty's intensity. Empirically, our results underscore the necessity of this weighting scheme; omitting the weighting coefficient—as in a simplistic variant—tends to induce degeneration in the language model’s output (see Appendix Table~\ref{tab:unlikelihood_generations}).

\textbf{DPO overview.} The standard DPO workflow consists of the following steps: First, for each prompt $x$, generate two completions $y_1, y_2 \sim \pi_\text{ref}(\cdot \mid x)$, annotate them according to human preferences, and form an offline preference dataset $\mathcal{D} = \{x^{(i)}, y_w^{(i)}, y_l^{(i)}\}_{i=1}^N$. Second, the language model $\pi_\theta$ is trained by minimizing $\mathcal{L}_\text{DPO}$ using the chosen $\pi_\text{ref}$, dataset $\mathcal{D}$, and a specified $\beta$. 

In practice, rather than collecting new completions and preference labels, existing public preference datasets are commonly utilized. These datasets typically originate from $\pi^\text{SFT}$; hence, when possible, we initialize $\pi_\text{ref}$ with $\pi^\text{SFT}$. When $\pi^\text{SFT}$ is unavailable, we initialize $\pi_\text{ref}$ by maximizing the likelihood of preferred completions ${(x, y_w)}$ over the dataset: ${\pi_\text{ref} = \argmax_{\pi}\mathbb{E}_{x, y_w \sim \mathcal{D}}\left[\log \pi(y_w \mid x)\right]}$. This strategy addresses discrepancies between the unknown true reference distribution and the actual $\pi_\text{ref}$ employed by DPO. Additional implementation details and hyperparameter choices are provided in Appendix~\ref{app:implementation}.

\section{Theoretical Analysis of DPO}
In this section, we further analyze the DPO approach, offering theoretical justification and discussing its strengths relative to the weaknesses of actor-critic algorithms applied in RLHF (for example, PPO~\cite{schulman2017proximal}).

\label{sec:theory}

\subsection{Your Language Model Is Secretly a Reward Model} 
DPO sidesteps both explicit reward model training and separate RL-based policy optimization, leveraging a unified maximum likelihood objective. The minimization objective described in Eq. \ref{eq:main_eq} can be seen as a Bradley-Terry model where the reward is defined as $r^*(x, y) = \beta \log\frac{\pi^*_\theta(y \mid x)}{\pi_\text{ref}(y \mid x)}$. Optimization of our parameterized model $\pi_{\theta}$ thus parallels the optimization of a reward model in Eq. \ref{eq:reward_model}, under an appropriate variable substitution. In the following, we develop the theoretical underpinnings of this reparameterization, demonstrate its lack of restrictiveness regarding the reward model family, and show that it admits exact retrieval of the optimal policy. We commence by introducing an equivalence relationship for reward functions.

\begin{definition}
Two reward functions $r(x, y)$ and $r'(x, y)$ are defined as equivalent if and only if ${r(x, y)-r'(x, y) = f(x)}$ for some function $f$.
\end{definition}

It is straightforward to verify that this definition constitutes an equivalence relation, thereby dividing the set of all reward functions into distinct equivalence classes. With this framework, we present the following two lemmas:

\begin{lemma}\label{lemma:same_prefrence}
Within the Plackett-Luce model, and specifically the Bradley-Terry setting, any pair of reward functions belonging to the same equivalence class induce identical preference distributions.
\end{lemma}

\begin{lemma}\label{lemma:same_policy}
    If two reward functions fall into the same equivalence class, then under the constrained RL formulation, they yield the same optimal policy.
\end{lemma}

Proofs for these statements are straightforward and provided in Appendix \ref{app:lemma1}. The first lemma highlights a well-recognized indeterminacy in the Plackett-Luce class of models \cite{plackett1975analysis}. This ambiguity means it is typically necessary to enforce extra identifiability constraints in order to obtain meaningful maximum likelihood estimates from Eq. \ref{eq:reward_model} \cite{bong2022generalized}. The second lemma shows that all members of a given equivalence class give rise to the same optimal policy; consequently, for our main objective, any representative reward function from this optimal class suffices. The following theorem, proven in Appendix~\ref{app:thm1}, formalizes this result:

\begin{theorem}\label{thm:main}
    Assuming mild regularity conditions, every equivalence class of reward functions consistent with Plackett-Luce (and especially Bradley-Terry) models can be expressed using the parameterization ${r(x, y) = \beta \log \frac{\pi(y\mid x)}{\pi_\text{ref}(y\mid x)}}$ for some policy model $\pi(y\mid x)$ and a fixed reference policy $\pi_\text{ref}(y \mid x)$.
\end{theorem}

\begin{sproof}
    Let $r(x, y)$ denote a reward function that determines a unique optimal policy $\pi_r(y \mid x)$ as prescribed by Eq. \ref{eq:op_policy}. We will demonstrate that a reward function from the equivalence class of $r$ can be recast via the aforementioned reparameterization. Specifically, we define the projection $f$ as  
\begin{equation}
    f(r; \pi_\text{ref}, \beta)(x, y) = r(x, y) - \beta\log\sum_{y}\pi_\text{ref}(y\mid x)\exp\left(\frac{1}{\beta}r(x, y)\right)
\end{equation}
This operator $f$ effectively adjusts the reward by subtracting the log of the partition function associated with $\pi_r$. Notably, since the normalization term depends only on $x$, $f(r; \pi_\text{ref}, \beta)(x, y)$ remains within the same equivalence class as $r(x, y)$. Substituting $r$ using the right-hand side of Eq.~\ref{eq:main_eq} (which applies to all reward functions) results in $f(r; \pi_\text{ref}, \beta)(x, y) = \beta \log \frac{\pi_r(y\mid x)}{\pi_\text{ref}(y\mid x)}$. Thus, the projection $f$ yields an element of the equivalence class for $r$ with the desired reparameterized form, ensuring the generality of our reward modeling approach under the proposed reparameterization.
\end{sproof}

An alternative interpretation of Theorem~\ref{thm:main} is that it identifies the specific reward function, from each equivalence class, that the DPO reparameterization picks. Concretely, it selects the reward function that fulfills the following condition:

\begin{equation}\label{eq:lag_p}
     \sum_{y}\underbrace{\pi_\text{ref}(y\mid x)\exp\left(\frac{1}{\beta}r(x, y)\right)}_{=\pi(y\mid x)\text{, as given by Thm.~\ref{thm:main} reparam.}} = 1,
\end{equation}

which ensures that $\pi(y\mid x)$ forms a legitimate probability distribution—that is, all probabilities are non-negative and collectively sum to one. Inspecting Eq.~\ref{eq:lag_p} through the lens of Eq.~\ref{eq:op_policy}, this equation is precisely the partition function associated with the optimal policy generated by the reward function $r(x, y)$. The principal contribution of the DPO approach is to enforce specific constraints within the otherwise under-determined Plackett-Luce family of preference models (including Bradley-Terry as a special case). These constraints do not restrict the reward models that can be represented; rather, they explicitly render the optimal policy of Eq. \ref{eq:op_policy} computationally manageable for any given prompt $x$.

\subsection{Instability of Actor-Critic Algorithms}
The framework described here also enables the identification of sources of instability commonly observed in standard actor-critic approaches to RLHF, such as PPO. Considering the RL fine-tuning step discussed in Section \ref{section:prelims}, we can relate our analysis to the control as inference perspective described by \cite{levine2018reinforcement} for the constrained RL problem in \ref{eq:RL}. Assuming a parameterized policy $\pi_{\theta}(y\mid x)$, the objective is to minimize $\mathbb{D}_{\text{KL}}[\pi_{\theta}(y|x) \mid \mid \pi^*(y\mid x)]$, where $\pi^*$ denotes the optimal policy described in Eq. \ref{eq:optimum_model} as determined by the reward function $r_{\phi}(y, x)$. Through elementary manipulations, this objective leads to the following maximization problem:

\begin{equation}\label{eq:AC}
    \max_{\pi_{\theta}}\mathbb{E}_{\pi_{\theta}(y\mid x)}\left[\underbrace{r_{\phi}(x, y) -\beta\log\sum_{y}\pi_\text{ref}(y\mid x)\exp\left(\frac{1}{\beta}r_{\phi}(x, y)\right)}_{f(r_{\phi}, \pi_\text{ref}, \beta)} - \underbrace{\beta\log\frac{\pi_{\theta}(y\mid x)}{\pi_\text{ref}(y\mid x)}}_{\text{KL}}\right]
\end{equation}

The objective under consideration matches that optimized in earlier research 
\citep{ziegler2020finetuning, stiennon2022learning, bai2022training, ouyang2022training}, where the reward function $r_{\phi}$ conforms to a DPO-equivalent form. Within this framework, the normalization component in $f(r_{\phi}, \pi_\text{ref}, \beta)$ can be viewed as representing the soft value function for the reference policy $\pi_\text{ref}$. Although this normalization term does not influence the location of the optimal solution, omitting it can increase the variance of the policy gradient estimator, which could undermine stability during training. It is possible to address this by introducing a value function estimator; however, optimizing such estimators poses its own challenges. Another approach adopted in prior literature involves normalizing the reward with respect to human completions, which is tantamount to using a single-sample Monte Carlo approximation of the normalization term. By contrast, the DPO reformulation results in a reward formulation that obviates the need for any baseline estimates.

\section{Experiments}
Here, we present empirical results to assess DPO’s effectiveness in learning policies directly from preference data. Our first line of investigation uses a controlled environment for text generation, wherein we evaluate DPO’s sample efficiency in trading off reward maximization and KL-divergence minimization against standard preference optimization methods, such as PPO. Subsequently, we scale our evaluation to larger neural architectures and more challenging RLHF applications—including summarization and dialogue tasks. Remarkably, we observe that DPO, with minimal hyperparameter adjustment, frequently matches or surpasses the performance of robust methods like RLHF with PPO, as well as the strategy of selecting the highest-rewarded trajectory out of $N$ samples. Prior to sharing our findings, we detail the experimental protocol; supplementary methodological specifics can be found in Appendix~\ref{app:exp_details}.

\textbf{Tasks.} Our study examines three distinct open-ended text generation tasks. Across all tasks, the models optimize a policy using a dataset of preferences $\mathcal{D}=\bigl\{x^{(i)}, y_w^{(i)}, y_l^{(i)}\bigr\}_{i=1}^N$. In the case of \textbf{controlled sentiment generation}, $x$ denotes the opening segment of a movie review sourced from the IMDb dataset \cite{maas-EtAl:2011:ACL-HLT2011}, with the objective for the policy to continue $x$ into a review $y$ that expresses a positive sentiment. To facilitate a systematic evaluation, this experiment generates preference pairs of model outputs by utilizing a pre-trained sentiment classifier, enforcing the criterion $p(\text{positive}\mid x,y_w)>p(\text{positive}\mid x,y_l)$. For SFT, the GPT-2-large model is fine-tuned on the training subset of the IMDB reviews dataset until performance stabilizes (implementation details in App~\ref{app:sentiment_details}). In the \textbf{summarization} task, $x$ represents a Reddit forum post, and the task is to create a concise summary $y$ that captures the main ideas. We adhere to established methodology by leveraging the Reddit TL;DR summarization dataset \citep{volske-etal-2017-tl} in combination with human preference data curated by \citeauthor{stiennon2022learning}. The SFT baseline used here is a model adapted from human-generated Reddit post summaries\footnote{\url{https://huggingface.co/CarperAI/openai_summarize_tldr_sft}}, employing the TRLX \citep{leandro_von_werra_2023_7790115} toolkit for RLHF. Human judgments for preferences originate from \citeauthor{stiennon2022learning}, who evaluated outputs from a similarly fine-tuned SFT model. Lastly, for \textbf{single-turn dialogue}, $x$ constitutes a user prompt—ranging from scientific inquiries to personal advice requests. The policy’s responsibility is to construct a reply $y$ that is both informative and engaging. Experiments utilize the Anthropic Helpful and Harmless dialogue dataset \citep{bai2022training}, which comprises 170,000 human–assistant conversational transcripts. Each conversation concludes with two assistant-generated responses from a large (but unspecified) language model, accompanied by human-labeled preferences indicating which response is favored. As there is no pre-existing SFT model for this dataset, we construct an SFT model by fine-tuning a publicly available language model exclusively on the preferred completions.

\textbf{Evaluation.} We employ two distinct methodologies for evaluation in our experiments. To systematically compare how well each algorithm optimizes the constrained reward maximization objective, we conduct a controlled sentiment generation task. Here, each algorithm’s performance is visualized by plotting its achieved reward against the KL-divergence from the reference policy. This analysis is feasible because we have direct access to the ground-truth reward signal via a sentiment classifier. Conversely, outside controlled environments where a ground-truth reward is unavailable, we assess each algorithm based on its \textit{win rate} in comparison to a baseline policy. For both the summarization and single-turn dialogue settings, we use GPT-4 as a stand-in for human judgment: it evaluates summary quality and dialogue helpfulness, respectively. The baseline for summarization is the corresponding reference summary from the test set, and for dialogue, it is the annotator-preferred response in the test data. While recent literature indicates that language models can outperform traditional metrics as automatic evaluators \citep{Chen2023ExploringTU}, we also conduct a human evaluation (detailed in Sec.~\ref{sec:human-judgments}) to validate GPT-4’s use as an evaluator. Our findings indicate that GPT-4’s assessments are highly correlated with human judgments, often matching or exceeding inter-annotator reliability.

\textbf{Methods.} Beyond DPO, our evaluation covers multiple established techniques for training language models to align with human preferences. As the simplest approach, we utilize zero-shot prompting with \textbf{GPT-J} \citep{gpt-j} for summarization, and 2-shot prompting with \textbf{Pythia-2.8B} \citep{biderman2023pythia} for dialogue. We also include the \textbf{SFT} model and the \textbf{Preferred-FT} variant, which is obtained through supervised finetuning using the selected completion $y_w$, with $y_w$ sourced from SFT for sentiment and summarization tasks, or from a general LM for single-turn dialogue. The \textbf{Unlikelihood} approach~\citep{welleck2019neural} is another method we investigate, where the objective is to maximize the probability assigned to $y_w$ while simultaneously minimizing the probability of $y_l$, optionally scaled by a tunable coefficient $\alpha\in[0,1]$ on the 'unlikelihood' component. Further, we include \textbf{PPO} \citep{schulman2017proximal} using a reward model trained from preference data, as well as \textbf{PPO-GT}, which directly utilizes the ground-truth reward in the controlled sentiment experiment. For PPO-GT, we test both an off-the-shelf implementation \cite{leandro_von_werra_2023_7790115} and a customized variant that normalizes the rewards and applies additional hyperparameter tuning to enhance outcomes; these modifications are similarly applied to standard PPO with learned rewards. Finally, we employ the \textbf{Best of $N$} baseline, which involves generating $N$ candidate responses with the SFT model (or Preferred-FT in dialogue) and selecting the highest-rewarding one according to a learned reward model. While this method can yield strong results, it is not scalable even for modest $N$ due to the requirement of sampling $N$ completions for every test instance.

\subsection{How well can DPO optimize the RLHF objective?}

The typical KL-constrained reward maximization objective in RLHF methodologies serves to manage a tradeoff: maximizing reward while simultaneously preventing the policy from diverging excessively from a predefined reference policy. Thus, algorithmic comparison necessitates jointly considering both the attained reward and the incurred KL divergence; obtaining a marginally superior reward with a substantially elevated KL is not always beneficial. In Figure~\ref{fig:frontier-tldr-main}, the relationship between reward and KL for various algorithms in the sentiment domain is illustrated. Each algorithm is subjected to multiple training sessions, varying the policy conservativeness hyperparameter per run (target KL $\in\{3,6,9,12\}$ for PPO; $\beta \in \{0.05,0.1,1,5\}$ for DPO; $\alpha\in\{0.05,0.1,0.5,1\}$ for unlikelihood; random seeds are used for preferred-FT). In total, this hyperparameter grid covers 22 training runs. At every 100-step interval throughout training, continuing until convergence, each policy is assessed on a fixed test prompt set. During evaluation, both the mean reward—measured according to the true reward function—and the average sequence-level KL\footnote{Defined as the aggregate of KL-divergences across all timesteps.} relative to the reference policy $\text{KL}\left(\pi\mid \mid \pi_\text{ref}\right)$ are computed. Our findings indicate that DPO attains the most effective frontier, achieving superior rewards without sacrificing KL efficiency. This outcome is remarkable for several reasons: although both DPO and PPO target the same optimization problem, DPO consistently yields a more favorable reward/KL balance, strictly outperforming PPO. Additionally, DPO establishes a superior reward-KL frontier compared to PPO, \emph{even when PPO benefits from direct access to ground truth rewards} (PPO-GT).

\subsection{Can DPO scale to real preference datasets?}
\label{sec:dpo-real-datasets}
We proceed to assess the capability of DPO for fine-tuning on tasks such as summarization and single-turn dialogue. In the context of summarization, standard automated metrics like ROUGE have been shown to have limited alignment with human judgment~\citep{stiennon2022learning}, and earlier research suggests that tuning language models using PPO informed by human preferences leads to improved summary quality. To evaluate different training approaches, we generate outputs on the test partition of the TL;DR summarization dataset and calculate the average rate at which these outputs outperform the reference completions in the test set. For each method, outputs are sampled at temperatures ranging from 0.0 to 1.0, with the corresponding win rates depicted in Figure~\ref{fig:frontier-tldr-main} (right). All models, including DPO, PPO, and Preferred-FT, are fine-tuned starting from the same GPT-J SFT checkpoint\footnote{\url{https://huggingface.co/CarperAI/openai_summarize_tldr_sft}}. Our results indicate that DPO achieves an approximate win rate of 61\% at temperature 0.0, outperforming PPO, which reaches about 57\% at its optimal temperature of 0.0. Additionally, DPO attains a higher peak win rate when compared with the best of $N$ approach. It is important to point out that we did not conduct an extensive search over the $\beta$ hyperparameter for DPO, so there may be further potential gains. We also observe that DPO demonstrates greater robustness to temperature variation relative to PPO, whose performance tends to regress toward that of the unmodified GPT-J model as temperature increases. Notably, Preferred-FT provides only marginal improvements over the base SFT model. Finally, in human preference assessments presented in Section~\ref{sec:human-judgments}, DPO-generated samples at a temperature of 0.25 were preferred 58\% of the time in comparison with PPO-generated samples at the same temperature.

For single-turn dialogue, we assess various approaches using the portion of the Anthropic HH dataset \citep{bai2022training} test split that involves only a single round of human-assistant exchange. To quantify win rates for each approach, we leverage GPT-4 comparisons with the reference being the preferred completions from the test set. Since a standard SFT baseline for this task does not exist, our experiments begin with a pre-trained Pythia-2.8B model. We then fine-tune a reference model using Preferred-FT on the selected completions to ensure they remain within the model's typical output distribution, followed by DPO training. Our analysis further includes a comparison to the top output among 128 Preferred-FT completions (having determined that the Best of $N$ approach reaches saturation at $N=128$ for this dataset; see Appendix Figure~\ref{fig:best-of-n}) and to a 2-shot prompt version of the Pythia-2.8B base model. Results indicate that DPO achieves equivalent or superior results compared to other methods when optimal temperature settings are chosen. Additionally, we consider an RLHF model optimized with PPO on the Anthropic HH dataset \footnote{\url{https://huggingface.co/reciprocate/ppo_hh_pythia-6B}} following a widely used implementation \footnote{\url{https://github.com/CarperAI/trlx/tree/main/examples/hh}}, but find that it does not outperform the base Pythia-2.8B under any prompt or temperature configuration. Given outcomes from TL;DR and the shared reward objective across methods, we use the Best of 128 setting as an approximate benchmark for PPO-level performance. In sum, DPO emerges as the only approach that not only surpasses the preferred completions in the Anthropic HH dataset efficiently, but also matches or exceeds the results of the far more computationally intensive Best of 128 baseline. Lastly, Figure~\ref{fig:dialogue-main} demonstrates that DPO reaches optimal performance after relatively few training iterations.

\subsection{Generalization to a new input distribution}

To extend our comparison between PPO and DPO in the presence of distributional changes, we assess both policies trained on Reddit TL;DR summarization by testing them on a different dataset, specifically, the test split of the CNN/DailyMail corpus \citep{nallapati-etal-2016-abstractive}. The evaluation uses the best sampling temperatures identified on the TL;DR set (0 and 0.25). Results are detailed in Table~\ref{tab:ood}. For evaluation, we used GPT-4 to determine win rates relative to reference summaries, applying the same GPT-4 (C) prompt as for Reddit TL;DR, but substituting “forum post” with “news article.” On this shifted input distribution, DPO remains consistently superior to the PPO policy by a notable margin. These findings suggest that DPO achieves generalization performance on par with, or better than, PPO—even though DPO does not rely on the supplementary unlabeled Reddit TL;DR prompts available to PPO.

\subsection{Validating GPT-4 judgments with human judgments}
\label{sec:human-judgments}

We implement a human evaluation study to assess the validity of GPT-4’s scoring, using outputs from the TL;DR summarization experiments along with two variations of GPT-4 prompts. The \textbf{GPT-4 (S)} (simple) prompt asks solely which summary better encapsulates the post’s key information. The \textbf{GPT-4 (C)} (concise) prompt extends this by additionally requesting judgment on conciseness. This modification is motivated by observations that, under the simple prompt, GPT-4 tends to favor lengthier, more repetitive summaries compared to human raters. Full prompt texts are available in Appendix~\ref{app:prompts}. Our study makes three comparisons to encompass a range of sample quality: the top-performing method (DPO, temp. 0.25), the lowest (PPO, temp. 1.0), and a mid-tier method (SFT, temp. 0.25), each evaluated against greedily-sampled PPO (chosen for its optimal temperature). The data reveal that, using both prompts, GPT-4’s judgments align with those of human raters as frequently as humans concur with one another, supporting the use of GPT-4 as a proxy for human evaluation (multiple human judgments are only collected for DPO and PPO-1 due to resource constraints). In summary, the \textbf{GPT-4 (C)} prompt yields win rates that more closely mirror human preferences, leading us to adopt it for principal reporting in Section~\ref{sec:dpo-real-datasets}. For more information on the human evaluation protocol, including rater interface and participant details, refer to Appendix~\ref{app:human-study}.

\section{Discussion}
Preference-based learning presents a robust and scalable approach for developing language models that are both competent and aligned. In this work, we have presented DPO, a streamlined training method that enables the learning of language models from preferences while forgoing reinforcement learning altogether. Instead of reframing the preference learning challenge to fit traditional RL frameworks and deploying standard RL algorithms, DPO establishes a direct correspondence between language model policies and reward representations. This correspondence allows for the direct optimization of language models in accordance with human preferences using a straightforward cross-entropy loss—eschewing reinforcement learning altogether and retaining full generality. Remarkably, DPO matches or surpasses existing RLHF methods, including those utilizing PPO, often requiring minimal hyperparameter tuning. Consequently, DPO significantly lowers the entry threshold for preference-driven language model training.

\textbf{Limitations \& Future Work.} Our findings prompt a number of compelling questions for subsequent research. For instance, how does a DPO-trained policy behave on out-of-distribution data in comparison with policies derived from explicit reward models? Preliminary evidence indicates DPO’s generalization may rival that of PPO-based approaches, but a more in-depth analysis is warranted. Another avenue concerns the potential of leveraging self-labeling with DPO policies to capitalize on unlabeled prompts efficiently. Further, there is an open question about how reward over-optimization emerges within the DPO paradigm—does the modest performance dip shown in Figure~\ref{fig:dialogue-main}-right reflect such an effect? Although our experiments have scaled up to models with 6B parameters, expanding DPO to much larger, state-of-the-art architectures remains a promising research trajectory. On the evaluation front, our observations indicate that GPT-4’s reported win rates are sensitive to the specifics of the evaluation prompt; future efforts could refine methodologies for eliciting robust automated assessments. Lastly, the potential applications of DPO are extensive and not limited to language modeling, as it may benefit the training of generative models across various modalities.